% Generated draft: requires mathematical and source review.
\documentclass[11pt]{article}
\usepackage[margin=25mm]{geometry}
\usepackage{fontspec}
\setmainfont{Latin Modern Roman}
\usepackage{amsmath,amssymb,mathtools,unicode-math}
\setmathfont{Latin Modern Math}
\usepackage{enumitem,needspace,xurl,hyperref}
\setlength{\parindent}{0pt}
\setlength{\parskip}{5pt}
\setlength{\emergencystretch}{3em}
\newcommand{\N}{\mathbb{N}}
\newcommand{\Z}{\mathbb{Z}}
\newcommand{\Q}{\mathbb{Q}}
\newcommand{\R}{\mathbb{R}}
\newcommand{\C}{\mathbb{C}}
\newcommand{\F}{\mathbb{F}}
\newcommand{\A}{\mathbb{A}}
\newcommand{\PP}{\mathbb{P}}
\newcommand{\E}{\mathbb{E}}
\newcommand{\eps}{\varepsilon}
\newcommand{\dd}{\,\mathrm{d}}
\begin{document}
\section*{UnsolvedMath ID: 30000681}
\textbf{Source problem code:} catalogue entry
\subsection{Definitions and mathematical statement}
\textbf{Mean-curvature locality question.} Let \(\mu_1,\mu_2\) be integral \(n\)-varifolds with generalized mean-curvature vectors \(H_{\mu_i}\in L^p(\mu_i)\). Is\[H_{\mu_1}(x)=H_{\mu_2}(x)\quad\text{for }\mu_1\text{-almost every }x\in\{\theta^{*n}(\mu_1)>0\}\cap\{\theta^{*n}(\mu_2)>0\}?\] The integrability exponent and ambient setting are not specified, and the source's one-sided almost-everywhere quantifier is retained.\par To clarify standard notation, take a common open ambient set \(U\subset\R^N\). An integral varifold \(V_i\) is an integer-multiplicity rectifiable \(n\)-varifold, and its weight measure \(\mu_i=\|V_i\|\) is defined on Borel sets \(E\subset U\) by \(\mu_i(E)=\int_{E\cap M_i}\vartheta_i\,d\mathcal H^n\), with \(M_i\) countably rectifiable and \(\vartheta_i\) positive integer-valued almost everywhere. The source uses \(\mu_i\) for both the varifold and the measure in density and integrability notation. Its first variation is \(\delta V_i(Y)=\int\operatorname{div}_T Y(x)\,dV_i(x,T)\) for compactly supported \(C^1\) vector fields. Under the absolutely continuous first-variation convention,\[\delta V_i(Y)=-\int_UH_{\mu_i}\cdot Y\,d\mu_i.\] If singular first variation is allowed, an additional singular term occurs and \(H_{\mu_i}\) denotes the absolutely continuous density instead. The record does not distinguish these conventions; no absence of singular variation is assumed silently.\par The star ordinarily denotes upper density, \(\theta^{*n}(\mu_i,x)=\limsup_{r\downarrow0}\mu_i(B_r(x))/(\omega_nr^n)\), where \(\omega_n\) is unit-ball volume. The source calls it \(n\)-density without fixing upper versus actual density, so the displayed limsup convention requires confirmation. In the usual rectifiable overlap, the positive-multiplicity measures on \(M_1\cap M_2\) are mutually absolutely continuous relative to \(\mathcal H^n\); equality of the vector fields is understood modulo the relevant null sets, not as pointwise equality of arbitrarily chosen representatives off their carriers. The question allows different multiplicities and does not assume identical varifolds in a neighborhood.\par The source gives no range for \(p\), codimension restrictions, first-variation hypotheses, or global versus localized interpretation beyond the stated \(L^p\) notation. The requested statement is locality of weak mean curvature on a measurable overlap, not merely the classical fact for two coinciding smooth patches. The supplied triage checked 2026-08-21 labels it open and reports no located complete resolution, with the report as sole reference. No precise later locality theorem is supplied; this dated status claim requires independent verification.
\subsection{Short English statement}
Ask whether the weak mean-curvature vectors of two integral varifolds agree almost everywhere on their positive-density overlap. The exponent and first-variation convention are missing. The comparison is between measure-defined vectors, not arbitrary pointwise representatives or only smooth coincident patches.
\subsection{Sources}
\begin{itemize}\item Local UnsolvedMath registry, permanent ID 30000681, code catalogue entry; \emph{Agreement of Mean Curvatures on Overlapping Varifolds}, Oberwolfach Reports Open Problems.\item \emph{Trends in Mathematical Imaging and Surface Processing (2007)}, Oberwolfach Reports, DOI 10.4171/owr/2007/03, pp. 139--194. \url{https://doi.org/10.4171/owr/2007/03}. Proposed year 2007; no proposer named.\item The same full report is the sole supplied literature source; no separate authoritative mean-curvature-locality theorem is supplied.\item Literature triage checked 2026-08-21, status \texttt{open}; record updated 2026-09-16. The \texttt{corrected\_verified} field claims editorial and notation cleanup only. Verify dated claims independently.\end{itemize}
\end{document}
